\section*{Capítulo \ref{ch:b2:himalaya-andes} --- Cordilleras de colisión y cordilleras andinas comparadas}

\begin{solution}{exo:b2:himalaya-andes:1}
(a) $r = h\rho_c/(\rho_m - \rho_c) = 3.0\times2700/600 = \qty{13.5}{km}$. (b) $T = 35 + 3.0 +
13.5 = \qty{51.5}{km}$, o bien $35 + 5.5\times3.0$. (c) $h = \Delta T(\rho_m - \rho_c)/\rho_m =
10\times600/3300 = \qty{1.8}{km}$.
\end{solution}

\begin{solution}{exo:b2:himalaya-andes:2}
(a) Antefosa indogangética; Cabalgamiento Frontal Principal; Subhimalaya (Siwalik); Cabalgamiento Principal
de Borde; Bajo Himalaya; Cabalgamiento Central Principal; secuencia del Alto Himalaya; Sistema de \omterm{def:b2:alpine-orogen:decollement}{Despegue}
del Tíbet Meridional; secuencia del Himalaya Tetisiano; después, la sutura del Indo--Tsangpo y la corteza
asiática. (b) El Sistema de \omterm{def:b2:alpine-orogen:decollement}{Despegue} del Tíbet Meridional. (c) El grado aumenta hacia el norte: \omterm{def:b2:alpine-orogen:foreland}{molasa} no metamorfizada,
Bajo Himalaya de grado medio, secuencia del Alto Himalaya de alto grado y parcialmente
fundida; por encima del STDS cae bruscamente a sedimentos tetisianos no metamorfizados o débilmente
metamorfizados, ya que la falla normal ha eliminado las rocas intermedias.
\end{solution}

\begin{solution}{exo:b2:himalaya-andes:3}
(a) $196.5\ \text{m}/9200\ \text{años} = \qty{0.02136}{m}$ al año, cerca de \qty{21.4}{mm} al año.
(b) Es un \qty{20}{\percent} más que los \qty{17.8}{mm} al año del GPS, pero del mismo orden:
una velocidad es un promedio sobre nueve mil años, la otra sobre veinte.
(c) Casi todo el acortamiento a través de la cordillera se absorbe en el frente: el MBT, el MCT y
las demás fallas al norte del MFT están, a esta longitud, prácticamente inactivas; todas enraízan
en el mismo MHT, que aflora en el MFT.
\end{solution}

\begin{solution}{exo:b2:himalaya-andes:4}
(a) Colisión: la sutura marca el océano desaparecido. (b) Tipo andino: un arco necesita una placa
que se hunda. (c) Colisión, típicamente: fundidos corticales de una corteza engrosada y calentada (también
aparecen en la corteza andina engrosada, pero allí dominan los \omterm{def:b2:magma-emplacement:magma}{magmas} de arco). (d) Tipo andino.
(e) Ambos: los Siwalik y el antepaís subandino. (f) Ambos: el MHT en el Himalaya, la
interfaz de subducción en los Andes.
\end{solution}

\begin{solution}{exo:b2:himalaya-andes:5}
(a) $17.8/45 = \qty{40}{\percent}$ en el este, $20.5/45 = \qty{46}{\percent}$ en el oeste:
cerca de la mitad. (b) Al engrosamiento y a las fallas del Tíbet y del resto de Asia, y a
la extrusión hacia el este de bloques a lo largo de las grandes fallas de desgarre. (c) $45 - 17.8 =
\qty{27.2}{mm}$ al año, o \qty{272}{km} en \qty{10}{Myr}; $45 - 20.5 = \qty{24.5}{mm}$ al
año, o \qty{245}{km}.
\end{solution}

\begin{solution}{exo:b2:himalaya-andes:6}
(a) $r = 4.5\times3.7 = \qty{16.65}{km}$; $\Delta F = 2700\times9.8\times3700\times(35\,000 +
1850 + 8325) = \qty{4.42e12}{N.m^{-1}}$. (b) Tíbet: \qty{6.12e12}{N.m^{-1}}; el Altiplano
tiene un \qty{72}{\percent} de ese valor. (c) Con $r = 4.5h$, $\Delta F = \rho_c g\,(Hh + 2.75h^2)$: para
$h \ll H$ domina el término $Hh$ y $\Delta F$ es proporcional a $h$; para una meseta alta
el término en $h^2$ crece más deprisa. Cada kilómetro adicional de altura cuesta cada vez más empuje en
los bordes, hasta que las fuerzas disponibles ya no pueden elevar la meseta; entonces se extiende
en lugar de elevarse.
\end{solution}

\begin{solution}{exo:b2:himalaya-andes:7}
(a) $2700\times9.8\times8848.86/2 = \qty{1.17e8}{Pa} = \qty{117}{MPa}$. (b) Con
\qty{100}{MPa}: $0.85\times100 = \qty{85}{MPa}$; con \qty{200}{MPa}: $0.85\times200 = 50 +
0.6\times200 = \qty{170}{MPa}$, las dos ramas coinciden. (c) $h = 2\sigma_d/(\rho g) =
2\times1.70\times10^8/(2700\times9.8) = \qty{12.8}{km}$: un orden de magnitud, del orden
de las cumbres más altas, no muy por encima de ellas.
\end{solution}

\begin{solution}{exo:b2:himalaya-andes:8}
(a) $11\times111.19\times\cos\ang{30} = \qty{1059}{km}$. (b) $9\times10^{-3}/1.059\times10^{6}
= \num{8.5e-9}$ al año (entre \num{6.6e-9} y \num{10.4e-9} para 7 y \qty{11}{mm} al año),
es decir, $8.5\times10^{-9}/3.156\times10^{7} = \qty{2.7e-16}{s^{-1}}$. (c) Los dos intervalos
se solapan: $(12.4 \pm 4.7) \times 10^{-9}$ va de 7.7 a \num{17.1e-9}. El sur de la meseta al norte de
Nepal se estira tan deprisa como el conjunto de la franja, o algo más.
\end{solution}

\begin{solution}{exo:b2:himalaya-andes:9}
(a) $0.9/6.6 = \qty{14}{\percent}$. (b) $M_0 = 10^{1.5\times8.5 + 9.1} = 10^{21.85} =
\qty{7.1e21}{N.m}$. (c) $7.1\times10^{21}/6.6\times10^{19} = 107$: un terremoto así cada
110 años aproximadamente bajo Nepal, muchos más de los que registra la historia; o bien se esperan terremotos mucho mayores,
o bien parte del déficit se libera de otro modo.
\end{solution}

\begin{solution}{exo:b2:himalaya-andes:10}
(a) Tómese la profundidad de compensación en la base de la corteza del Altiplano, $61.5 - 3.7 =
\qty{57.8}{km}$ bajo el nivel del mar. La corteza de la Puna llega a $50.5 - 4.4 = \qty{46.1}{km}$, de modo que la
columna de la Puna contiene \qty{11.7}{km} de manto por encima de esa profundidad: $2800\times61.5 =
\rho_P\times50.5 + 3300\times11.7$. (b) $\rho_P = (172\,200 - 38\,610)/50.5 =
\qty{2645}{kg.m^{-3}}$. (c) La diferencia, \qty{155}{kg.m^{-3}}, está en el extremo inferior de
los 150 a \qty{200}{kg.m^{-3}} citados: la corteza sola requeriría una corteza de la Puna inusualmente ligera.
Una columna más ligera bajo ella lo consigue con más facilidad: un manto litosférico eliminado por delaminación,
sustituido por astenosfera caliente y ligera.
\end{solution}

\begin{solution}{exo:b2:himalaya-andes:11}
(a) $g = 105/\tan\ang{30} = \qty{181.9}{km}$. (b) $g = 90/\tan\ang{30} + 15/\tan\ang{4} =
155.9 + 214.5 = \qty{370.4}{km}$: el arco se desplaza \qty{188.5}{km} hacia el interior. (c) La placa nunca
alcanza los \qty{105}{km} bajo la placa superior a lo largo de esos \qty{300}{km}, y no hay ninguna cuña de manto caliente
encima de ella: no se produce \omterm{def:b2:magma-emplacement:magma}{magma} y el arco se extingue hasta que la placa vuelve a hundirse, más al este.
\end{solution}

\begin{solution}{exo:b2:himalaya-andes:12}
(a) \qty{21}{mm} al año son \qty{21}{km} por millón de años: $21\times37.5 =
\qty{787.5}{km^3}$ por kilómetro de cordillera y por millón de años. (b) $300/787.5 =
\qty{38}{\percent}$. (c) En el engrosamiento de la cordillera y del sur del Tíbet, en la corteza
arrastrada con la India bajo el Tíbet, y en el material que el MHT transfiere hacia el norte.
\end{solution}

\begin{solution}{pb:b2:himalaya-andes:1}
\textbf{1.} Siete muestras de \ang{28.0} a \ang{37.0}~N: $34\,550/7 = \qty{4936}{m}$, $h_T
= \qty{4.9}{km}$.
\textbf{2.} Cinco muestras de \ang{-69} a \ang{-65}: $18\,650/5 = \qty{3730}{m}$, $h_A =
\qty{3.7}{km}$, el valor citado para el Altiplano; la Puna, más al sur, es más alta.
\textbf{3.} Tíbet: al menos $9\times111.19 = \qty{1001}{km}$; Altiplano: al menos $4\times
104.5 = \qty{418}{km}$.
\textbf{4.} $(5000 - 70)/(1.5\times111\,190) = 0.0296$: un \qty{3.0}{\percent}.
\textbf{5.} El perfil andino empieza en el fondo oceánico y en la fosa, kilómetros por debajo
del nivel del mar, al oeste de la meseta; el tibetano empieza en una llanura continental cerca del nivel del mar,
y termina en la alta cuenca del Tarim.
\textbf{6.} $r = 4.5\times4.9 = \qty{22.05}{km}$; $T_T = 35 + 4.9 + 22.05 = \qty{62.0}{km}$.
\textbf{7.} $r = 4.5\times3.7 = \qty{16.65}{km}$; $T_A = 35 + 3.7 + 16.65 = \qty{55.4}{km}$.
\textbf{8.} Tíbet: \qty{62.0}{km} frente a 70--\qty{75}{km}, \qty{10.5}{km} por debajo del valor
medio \qty{72.5}{km} (un \qty{14}{\percent}); Altiplano: \qty{55.4}{km} frente a 59--\qty{64}{km},
\qty{6.1}{km} por debajo de \qty{61.5}{km} (un \qty{10}{\percent}). El Tíbet se aleja más.
\textbf{9.} $72.5 - 35 = 4.9\times3300/(3300 - \rho_c)$: $3300 - \rho_c = 16\,170/37.5 =
431.2$, $\rho_c = \qty{2869}{kg.m^{-3}}$.
\textbf{10.} $26.5 = 3.7\times3300/(3300 - \rho_c)$: $3300 - \rho_c = 460.8$, $\rho_c =
\qty{2839}{kg.m^{-3}}$.
\textbf{11.} Parte de la carga la soporta la flexión de la placa, no una raíz local; la
corteza normal puede no haber tenido \qty{35}{km} de espesor; el manto de debajo puede ser más ligero o más denso
de lo supuesto; el Moho sísmico puede no ser la base de la corteza en el sentido de la densidad.
\textbf{12.} La base de la corteza del Tíbet está a $72.5 - 4.9 = \qty{67.6}{km}$, y la de la India a
\qty{37.5}{km}: $t = \qty{30.1}{km}$. $\Delta g = 2\pi\times6.674\times10^{-11}\times370\times
30\,100 = \qty{4.67e-3}{m.s^{-2}}$, es decir, \qty{-467}{mGal} respecto de la India, cerca de los
\qty{460}{mGal} observados.
\textbf{13.} $W_0 = 440\times60/35 = \qty{754.3}{km}$; $S = \qty{314.3}{km}$.
\textbf{14.} $(314.3 - 250)/314.3 = \qty{20}{\percent}$, el déficit hallado en los Andes
Centrales.
\textbf{15.} $0.55\times314.3 = \qty{173}{km}$, menos que los \qty{250}{km} de los \omterm{def:b2:alpine-orogen:balanced}{cortes
compensados}, que serían entonces más que suficientes.
\textbf{16.} La cobertera acortada por encima del \omterm{def:b2:alpine-orogen:decollement}{despegue} debe compensarse con un acortamiento del
zócalo de debajo, más al oeste, ya que ambos deben perder la misma longitud a través de la cordillera; el acortamiento
del zócalo engruesa la corteza bajo la Cordillera Oriental y la meseta.
\textbf{17.} $314.3/30 = \qty{10.5}{mm}$ al año, $10.5/72 = \qty{15}{\percent}$ de la
convergencia.
\textbf{18.} Se consume por subducción: la placa de Nazca se hunde bajo el continente.
\textbf{19.} $r = \qty{22.05}{km}$; $\Delta F = 2700\times9.8\times4900\times(35\,000 + 2450
+ 11\,025) = \qty{6.28e12}{N.m^{-1}}$.
\textbf{20.} $2700\times9.8\times3700\times45\,175 = \qty{4.42e12}{N.m^{-1}}$; la razón es
0.70.
\textbf{21.} $6.28\times10^{12}/10^{5} = \qty{6.3e7}{Pa}$, unos \qty{63}{MPa}.
\textbf{22.} $9\times10^{-3}/1.06\times10^{6} = \num{8.5e-9}$ al año.
\textbf{23.} El empuje de la India actúa en dirección norte--sur y se equilibra, en la alta meseta, con el
propio peso de la meseta: en dirección norte--sur la meseta ya no puede engrosarse, de modo que el esfuerzo vertical,
el peso, es el mayor; el esfuerzo menor es este--oeste, hacia los bordes libres de Asia,
y la meseta se extiende en esa dirección, mediante rifts norte--sur.
\textbf{24.} \textbf{\qty[text-series-to-math=true]{62}{km}}, cerca de un \qty{86}{\percent} de los
\qty{72.5}{km} medidos.
\end{solution}
